% Построение линейной цепи роя по двум якорям (сценарий linear_chain2).
% Сборка: pdflatex linear_chain_two_anchors.tex
\documentclass[12pt,a4paper]{article}

\usepackage[utf8]{inputenc}
\usepackage[T2A]{fontenc}
\usepackage[russian]{babel}
\usepackage{amsmath,amssymb,amsthm}
\usepackage{geometry}
\geometry{margin=2.5cm}
\usepackage{hyperref}
\hypersetup{colorlinks=true,linkcolor=blue,urlcolor=blue}

\title{Алгоритм выстраивания роя в линию по двум якорям\\
(геометрия отрезка и закон управления в симуляторе)}
\author{По коду \texttt{Drone\_Swarm\_Simulator\_v2/scenarios/linear\_chain2.py}}
\date{}

\begin{document}
\maketitle

\begin{abstract}
Рассматривается упрощённая постановка: два концевых борта задают концы желаемой линии; остальные дроны выравниваются на прямую, проходящую через эти точки, и стремятся к равномерному распределению вдоль отрезка между ними. Дополнительно используется локальная нелинейная поправка по ближайшим видимым соседям в базисе «вдоль/поперёк'' цепи --- как вспомогательный стабилизирующий вклад, а не как описание обхода коридора. Ниже зафиксированы обозначения, геометрия, целевые точки якорей и итоговый вектор команды в общей горизонтальной плоскости NED с соответствием именам в исходном коде.
\end{abstract}

%---------------------------------------------------------------------
\section{Постановка задачи и обозначения}

Пусть в эксперименте участвует $N \ge 4$ однотипных мультикоптеров. Два из них играют роль \emph{якорей-концов} цепи: в коде это борты с минимальным и максимальным идентификатором в конфигурации (\texttt{endpoint\_left\_id}, \texttt{endpoint\_right\_id}). Остальные --- \emph{внутренние} агенты.

Горизонтальные координаты для логики цепи задаются в \emph{общей} системе NED плоскости $Oxy$: ось $x$ --- север, ось $y$ --- восток (как в \texttt{LOCAL\_POSITION\_NED}). Для согласования визуального разнесения бортов в симуляторе к координате $y$ каждого дрона добавляется поправка, зависящая только от его номера (\texttt{\_did\_offset\_y}); далее под $p_i = (x_i, y_i)^\top$ понимается уже скорректированное положение в этой общей плоскости.

Состояние концевых дронов в каждый момент принимается за текущие положения концов желаемого отрезка. Обозначим $A(t)$ --- положение левого якоря, $B(t)$ --- правого (в общих координатах). Внутренние агенты не задают глобальную траекторию коридора; их задача --- согласоваться с прямой $AB$ и промежуточными положениями вдоль неё.

В теоретических работах о саморазвертывании между двумя границами часто фигурируют расстояния до ближайших соседей слева и справа и монотонная функция от этих расстояний, обеспечивающая в целевом режиме почти равномерный шаг на интервале между опорными точками. Реализация в \texttt{linear\_chain2} опирается на ту же идею \emph{локальных} соседей, но основной упор в управлении делается на явную геометрию отрезка $AB$.

%---------------------------------------------------------------------
\section{Геометрия двух якорей: отрезок $AB$, базис $(\mathbf{u}, \mathbf{n})$, проекция $s_i$}

Единичный вектор вдоль цепи от левого якоря к правому задаётся нормировкой разности положений:
\begin{equation}
  \mathbf{u} = \frac{B - A}{\lVert B - A \rVert}, \qquad
  \mathbf{u} \in \mathbb{R}^2, \quad \lVert \mathbf{u} \rVert = 1.
\end{equation}
Вектор $\mathbf{u} = (u_x, u_y)^\top$ строится в коде как \texttt{\_unit\_vector\_xy(pos\_l, pos\_r)}.

Перпендикуляр к $\mathbf{u}$ в плоскости $Oxy$, ориентированный «влево'' от направления $\mathbf{u}$ (поворот на $+90^\circ$ в плоскости экрана север--восток), задаётся
\begin{equation}
  \mathbf{n} = (-u_y,\, u_x)^\top
\end{equation}
(функция \texttt{\_normal\_left\_xy}).

Проекция положения $i$-го дрона на ось, проходящую через точку $A$ в направлении $\mathbf{u}$, определяется скаляром
\begin{equation}
  s_i = (p_i - A) \cdot \mathbf{u}.
\end{equation}
В коде это \texttt{\_projection\_s(endpoint\_a, u, p)} с \texttt{endpoint\_a = pos\_l}.

Поперечное (подписанное) отклонение от прямой $AB$ удобно записать как
\begin{equation}
  e_{\mathrm{cross},\,i} = (p_i - A) \cdot \mathbf{n};
\end{equation}
оно совпадает со знаковым расстоянием до прямой, проходящей через $A$ вдоль $\mathbf{u}$, если $\lVert \mathbf{n} \rVert = 1$.

%---------------------------------------------------------------------
\section{Движение якорей и целевые точки в симуляторе}

После инициализации и обмена координатами для концевых бортов вычисляются \emph{целевые} точки $A^\star$, $B^\star$ в той же общей плоскости (\texttt{\_anchor\_targets\_from\_initial\_segment}). Пусть $L$ --- параметр длины сегмента \texttt{segment\_length} (в метрах); $L_{\mathrm{half}} = L/2$. Если в момент расчёта якоря находятся в скорректированных точках с координатами $(x_A, y_A)$ и $(x_B, y_B)$, то
\begin{align}
  A^\star &= \bigl(x_A,\; y_A - L_{\mathrm{half}},\; z_\star\bigr)^\top, \\
  B^\star &= \bigl(x_B + L_{\mathrm{half}},\; y_B,\; z_\star\bigr)^\top,
\end{align}
где $z_\star$ --- заданная высота удержания (в коде соответствует отрицательной высоте NED после взлёта).

Каждый якорный борт переводится к своей цели с помощью ПИД-закона по ошибке положения в горизонтальной плоскости: сначала ошибка $(e_x, e_y)^\top = (x^\star - x,\, y^\star - y)^\top$ в NED преобразуется в «вперёд--вправо'' по текущему рысканию (\texttt{\_ned\_to\_body\_forward\_right}), затем по осям корпуса вычисляются выходы двух регуляторов (\texttt{endpoint\_loop}, \texttt{\_move\_towards\_common\_with\_pid}). После достижения окрестности цели (порог \texttt{ENDPOINT\_REACHED\_THRESH\_M}) команды сбрасываются в нейтраль. Поэтому корректно говорить, что концы цепи \emph{подводятся} к $A^\star$ и $B^\star$ и далее задают текущую геометрию отрезка $AB$ для внутренних агентов; активное ПИД-удержание цели после срабатывания порога в этом цикле уже не выполняется.

%---------------------------------------------------------------------
\section{Закон управления внутренних дронов}

Внутренний цикл (\texttt{internal\_chain\_anatoliy\_loop}) с частотой \texttt{CONTROL\_HZ} на каждом шаге заново вычисляет $\mathbf{u}$, $\mathbf{n}$, $A$ по текущим положениям концевых бортов, собирает множество «видимых'' соседей в круге радиуса $R_{\mathrm{vis}}$ (\texttt{r\_vis}) и формирует горизонтальный вектор комбинированного ускорения/направления $\mathbf{c}$ в NED, который затем поворачивается в корпус, демпфируется по измеренной скорости и преобразуется в \texttt{RC\_OVERRIDE}. Если положения якорей ещё не получены или в пределах видимости оказалось меньше двух соседей, борт получает нейтральные команды и ждёт следующего шага обновления.

Ниже $k_{\mathrm{along}} = \texttt{SPACING\_ALONG\_K}$, $k_{\mathrm{cross}} = \texttt{SPACING\_CROSS\_K}$, $w_A = \texttt{ANATOLIY\_COMB\_WEIGHT}$, $w_{\mathrm{cross}} = \texttt{W\_CROSS\_M}$.

\subsection{Равномерное распределение вдоль $AB$}

Все борты с известными позициями упорядочивают по возрастанию $s_i$ (при равенстве $s$ --- по номеру). Для внутреннего дрона $i$, у которого есть и левый, и правый сосед в этом порядке, вводится продольная ошибка равномерности
\begin{equation}
  e_{\mathrm{along},\,i} = \frac{s_{i^-} + s_{i^+}}{2} - s_i,
\end{equation}
где $i^-$, $i^+$ --- соседние в отсортированном списке идентификаторы, $s_{i^\pm}$ --- их проекции. Если соседа с одной из сторон нет, в реализации $e_{\mathrm{along},\,i} = 0$ (цепочка не дополняется искусственными границами по $s$ для крайних внутренних в особом виде --- крайние якоря задают геометрию через $A,B$). Величина $e_{\mathrm{along},\,i}$ ограничивается по модулю константой \texttt{SPACING\_E\_ALONG\_CLAMP\_M}.

Вклад в горизонтальную команду, отвечающий за выравнивание шага вдоль цепи:
\begin{equation}
  \mathbf{g}_{\mathrm{geom}} = k_{\mathrm{along}}\, e_{\mathrm{along},\,i}\, \mathbf{u} - k_{\mathrm{cross}}\, e_{\mathrm{cross},\,i}\, \mathbf{n}.
\end{equation}
Первое слагаемое тянет проекцию $s_i$ к середине между соседями по $s$, что в целевом режиме соответствует выравниванию интервалов между соседними дронами вдоль $AB$.

\subsection{Подавление отклонения от прямой $AB$}

Второе слагаемое в $\mathbf{g}_{\mathrm{geom}}$ использует $e_{\mathrm{cross},\,i}$: оно направлено вдоль $-\mathbf{n}$ с коэффициентом $k_{\mathrm{cross}}$, то есть уменьшает поперечное смещение относительно прямой через $A$ вдоль $\mathbf{u}$. Величина $e_{\mathrm{cross},\,i}$ также ограничивается по модулю (\texttt{SPACING\_E\_CROSS\_CLAMP\_M}).

\subsection{Вспомогательная локальная поправка по видимым соседям}

Для каждого другого дрона $j$ в пределах видимости вводятся относительные координаты в базисе $(\mathbf{u}, \mathbf{n})$:
\[
  \Delta p_{ij} = p_j - p_i, \quad
  a_{ij} = \Delta p_{ij} \cdot \mathbf{u}, \quad
  c_{ij} = \Delta p_{ij} \cdot \mathbf{n}.
\]
По множеству $\{(a_{ij}, c_{ij})\}$ выбираются четыре расстояния до ближайшего соседа в каждом из четырёх «квадрантов'' знаков $(a,c)$, аналогично разложению вдоль/поперёк в сценарии коридора: $d_a^+$, $d_c^+$, $d_a^-$, $d_c^-$ (в коде \texttt{\_nearest\_along\_cross}). Для каждой величины хранится флаг наличия соседа в соответствующем направлении ($f_a^\pm$, $f_c^\pm$).

Функция подстановки «расстояние или зазор'' (как $g$ в \texttt{corridor\_sweeping}) задаётся
\begin{equation}
  g(d, w, f) =
  \begin{cases}
    d, & f = \text{истина (сосед найден)}, \\
    w, & f = \text{ложь}.
  \end{cases}
\end{equation}
Для продольной оси в реализации при отсутствии соседа подставляется $w = 0$; для поперечной --- $w = w_{\mathrm{cross}}$.

Пусть $\xi(\cdot) = \tanh(\cdot)$. Скоростные члены в законе коридора в симуляторе положены в нуль (\texttt{v\_along = v\_cross = 0}), чтобы избежать дрейфа сигнала от шумной телеметрии скорости. Тогда вспомогательные скаляры принимают вид
\begin{align}
  \sigma_a &= -\xi\!\bigl(g(d_a^+, 0, f_a^+)\bigr) + \xi\!\bigl(g(d_a^-, 0, f_a^-)\bigr), \\
  \sigma_c &= -\xi\!\bigl(g(d_c^+, w_{\mathrm{cross}}, f_c^+)\bigr) + \xi\!\bigl(g(d_c^-, w_{\mathrm{cross}}, f_c^-)\bigr).
\end{align}
Вектор вспомогательной поправки в плоскости:
\begin{equation}
  \mathbf{a} = \sigma_a\, \mathbf{u} + \sigma_c\, \mathbf{n}.
\end{equation}

Итоговый горизонтальный вектор направления до поворота в корпус и демпфирования:
\begin{equation}
  \mathbf{c} = \mathbf{g}_{\mathrm{geom}} + w_A\, \mathbf{a}.
\end{equation}

Далее $\mathbf{c}$ переводится в оси «вперёд--вправо'' по рысканию (\texttt{\_ned\_to\_body\_forward\_right}), из него вычитается демпфирование, пропорциональное проекции скорости в тех же осях (\texttt{INTERNAL\_BODY\_V\_DAMP}), ошибки по осям масштабируются в приращения PWM (\texttt{INTERNAL\_RC\_KP}, ограничения и \texttt{\_slew\_int}), и отправляется \texttt{send\_rc\_override}.

Если и геометрическая ошибка мала (мёртвые зоны \texttt{SPACING\_E\_DEADBAND\_M}), и $|\sigma_a|$, $|\sigma_c|$ ниже \texttt{SIGMA\_DEADBAND}, борт получает нейтральные стики. Утверждения о строгой сходимости к равновесию здесь не формулируются: корректно говорить о \emph{целевом} выравнивании вдоль $AB$ и о роли $\mathbf{a}$ как дополнительной локальной нелинейной подстройки плотности.

%---------------------------------------------------------------------
\section{Соответствие формул строкам в \texttt{linear\_chain2.py}}

\begin{itemize}
  \item Общие координаты и сдвиг по $y$: \texttt{\_my\_position\_common}, \texttt{\_did\_offset\_y}; обмен --- \texttt{coordinate\_exchange\_loop}.
  \item $\mathbf{u}$, $\mathbf{n}$, $s_i$, $e_{\mathrm{cross}}$: \texttt{\_unit\_vector\_xy}, \texttt{\_normal\_left\_xy}, \texttt{\_projection\_s}, блок \texttt{\_spacing\_errors\_projection}.
  \item $e_{\mathrm{along}}$: сортировка \texttt{ids\_sorted} по $s$, формула \texttt{0.5 * (s\_lid + s\_rid) - s\_me}.
  \item Цели якорей $A^\star$, $B^\star$: \texttt{\_anchor\_targets\_from\_initial\_segment}; ПИД концов --- \texttt{endpoint\_loop}, \texttt{\_move\_towards\_common\_with\_pid}.
  \item $d_a^\pm$, $d_c^\pm$, флаги: \texttt{\_nearest\_along\_cross}; $g$, $\xi$: \texttt{\_g}, \texttt{\_xi}; $\sigma_a$, $\sigma_c$: \texttt{\_sigmas\_segment\_frame}.
  \item $\mathbf{g}_{\mathrm{geom}}$, $\mathbf{a}$, $\mathbf{c}$: переменные \texttt{gx, gy}, \texttt{ax, ay}, \texttt{comb\_x, comb\_y} с весами \texttt{SPACING\_*}, \texttt{ANATOLIY\_COMB\_WEIGHT}.
  \item Корпус и PWM: \texttt{\_yaw\_rad}, \texttt{\_ned\_to\_body\_forward\_right}, \texttt{\_velocity\_xy\_ned}, \texttt{\_pwm\_axis}, \texttt{\_slew\_int}, \texttt{send\_rc\_override}.
\end{itemize}

%---------------------------------------------------------------------
\section*{Простыми словами (для устного сопровождения)}

Два крайних дрона --- это «гвозди'': они уезжают на заранее заданные точки и задают, где должна проходить линия. Все остальные смотрят на эту линию как на дорожку: каждый старается встать ровно посередине между соседями вдоль дорожки и не уезжать вбок от неё. Дополнительно каждый учитывает, кто стоит рядом в поле зрения, через мягкую нелинейную функцию --- это мелкая подправка, чтобы не толкаться и не слипаться, а не отдельная «логика коридора''. Команда в итоге переводится в привычные для пилота наклоны коптера; скорость слегка гасится, чтобы не раскачивать звено.

\end{document}
